\documentclass[border=5mm]{standalone}
% ==============================================================================
% SETUP.TEX - CONFIGURACIÓN GLOBAL DEL PROYECTO
% ==============================================================================
% Este archivo centraliza todos los paquetes y estilos.
% Debe incluirse en el preámbulo del libro principal y de los archivos standalone.
% \PassOptionsToPackage{gray}{xcolor}
% --- 1. CODIFICACIÓN E IDIOMA ---
% fix-cm habilita las fuentes Computer Modern en tamaños arbitrarios (por debajo
% de 5 pt, entre otros). Debe cargarse antes que fontenc.
\usepackage{fix-cm}
\usepackage[utf8]{inputenc}
\usepackage[T1]{fontenc}
\usepackage[spanish, es-tabla]{babel}  
\usepackage{csquotes} % Recomendado para babel y citas

\usepackage{import}
% mode=image: Usa el PDF pre-compilado, nunca intenta recompilar.
% Debes compilar cada figura manualmente antes de compilar el libro.
\usepackage[mode=image, subpreambles=false]{standalone}

% --- 2. MATEMÁTICAS Y CIENCIAS ---
\usepackage{amsmath, amssymb, amsfonts, mathtools}
\usepackage{enumitem}

% LaTeX trae \DeclareMathSizes{5}{5}{5}{5}, así que a 5 pt (\tiny) los subíndices
% salen del mismo tamaño que la base: en las etiquetas de figuras el M de
% $\omega_M$ queda desproporcionado. Se restablece la proporción habitual.
\DeclareMathSizes{5}{5}{3.7}{3}

% --- 3. DISEÑO Y GEOMETRÍA --- 
% Unificamos los márgenes para todo el documento
\usepackage[left=2.5cm,right=2.5cm,top=3cm,bottom=3cm]{geometry}
\makeatletter
\ifstandalone % Evita que geometry dañe el recorte de las figuras bajándolas 3cm
  \@ifclasswith{standalone}{crop=false}{
    % Es un capítulo/sección: Dejamos que geometry actúe.
  }{
    % Es una figura: Neutralizamos geometry para que no las recorte mal.
    \geometry{pass}
  }
\fi
\makeatother
\usepackage{parskip} % Espaciado entre párrafos sin sangría
\usepackage{float}   % Para usar [H] en figura

% Ajustes de flotantes para evitar espacios verticales exagerados
\renewcommand{\topfraction}{0.95}      % Permite que las figuras ocupen hasta el 95% superior
\renewcommand{\bottomfraction}{0.95}   % Permite que las figuras ocupen hasta el 95% inferior
\renewcommand{\textfraction}{0.05}     % Permite hojas con tan solo un 5% de texto
\renewcommand{\floatpagefraction}{0.8} % Requiere que una página de solo figuras esté 80% llena
\setcounter{topnumber}{4}
\setcounter{bottomnumber}{4}
\setcounter{totalnumber}{10}

% --- 4. GRÁFICOS Y TABLAS ---
\usepackage{graphicx}
\usepackage{booktabs} % Tablas profesionales
\usepackage{caption}  % Control de captions y \captionof
\usepackage{tikz}
\usetikzlibrary{arrows.meta, calc, angles, quotes, babel, backgrounds}
\usepackage{pgfplots}
\usepgfplotslibrary{groupplots}
\usepgfplotslibrary{external}
\usepgfplotslibrary{patchplots}
\pgfplotsset{compat=1.18}
\usepackage[american, straightvoltages]{circuitikz}

% --- 5. COLORES Y CAJAS ---

\usepackage[table]{xcolor}
\usepackage{tcolorbox}

% Definición de colores corporativos/estilo
\definecolor{utnblue}{RGB}{0, 51, 102}
\definecolor{codegreen}{rgb}{0,0.6,0}
\definecolor{codegray}{rgb}{0.5,0.5,0.5}
\definecolor{codepurple}{rgb}{0.58,0,0.82}
\definecolor{backcolour}{rgb}{0.96,0.96,0.96}

% --- 6. ENTORNOS PERSONALIZADOS (CAJAS) ---

% Caja "Resultado" (Azul Corporativo) — Definiciones, fórmulas clave, resultados importantes.
% Uso: \begin{resultado}[Fórmula de Euler] ... \end{resultado}
\newtcolorbox{resultado}[1][]{
    colback=utnblue!5!white,
    colframe=utnblue,
    title=\textbf{#1},
    fonttitle=\bfseries,
    sharp corners,
    boxrule=0.5mm
}

% Caja "Nota" (Gris) — Advertencias, aplicaciones, contexto secundario.
% Uso: \begin{nota}[Cuidado con la calculadora] ... \end{nota}
% Uso: \begin{nota} ... \end{nota}  → título por defecto "Nota"
\newtcolorbox{nota}[1][Nota]{
    colback=codegray!5!white,
    colframe=codegray,
    title=\textbf{#1},
    fonttitle=\bfseries,
    sharp corners,
    boxrule=0.5mm
}

% Caja inline para resaltar un valor y enlazarlo con su otra aparición en una tabla
% o en una cuenta: mismo color, mismo valor.
% Uso: \marcavalor{$4$}  o  \marcavalor[codegreen]{$4$}
\newtcbox{\marcavalor}[1][utnblue]{
    on line,
    boxsep=1pt, left=2pt, right=2pt, top=1pt, bottom=1pt,
    colback=#1!10!white,
    colframe=#1,
    boxrule=0.4pt,
    arc=2pt
}

% --- 7. CÓDIGO FUENTE (LISTINGS) ---
\usepackage{listings}

\lstdefinestyle{miestilo}{
    backgroundcolor=\color{backcolour},   
    commentstyle=\color{codegreen},
    keywordstyle=\color{magenta},
    numberstyle=\tiny\color{codegray},
    stringstyle=\color{codepurple},
    basicstyle=\ttfamily\footnotesize,
    breakatwhitespace=false,         
    breaklines=true,                 
    captionpos=b,                    
    keepspaces=true,                 
    numbers=left,                    
    numbersep=5pt,                  
    showspaces=false,                
    showstringspaces=false,
    showtabs=false,                  
    tabsize=2,
    frame=single,                    
    rulecolor=\color{codegray}
}

\lstset{style=miestilo}
% Soporte para tildes y ñ en bloques de código
\lstset{literate={ñ}{{\~n}}1 {á}{{\'a}}1 {é}{{\'e}}1 {í}{{\'i}}1 {ó}{{\'o}}1 {ú}{{\'u}}1}

% Operadores matemáticos personalizados

% Vector en sentido discreto (lista finita de coordenadas): negrita + flecha.
% Uso: \vect{v}, \vect{e}_k, \vect{0}.
\newcommand{\vect}[1]{\vec{\mathbf{#1}}}

% Fecha de versión y comando para capítulos standalone
\providecommand{\verdate}{\the\year.\ifnum\month<10 0\fi\the\month.\ifnum\day<10 0\fi\the\day}
\newcommand{\chapterversion}{%
    \vspace{-1.5cm}\begin{flushright}\small\textit{\color{codegray}Versión~\verdate}\end{flushright}\vspace{0.5cm}%
}

% --- 8. HIPERVÍNCULOS (DEBE SER EL ÚLTIMO PAQUETE) ---
\usepackage[colorlinks=true, linkcolor=utnblue, urlcolor=utnblue, citecolor=utnblue]{hyperref}

% --- 9. REFERENCIAS CRUZADAS ENTRE CAPÍTULOS STANDALONE ---
% Cuando se compila un capítulo standalone, xr-hyper lee los .aux de los
% otros capítulos (si existen) para resolver \ref{} a labels externos.
% En la compilación del libro completo este bloque no se activa.
\ifstandalone
  \usepackage{xr-hyper}
  \makeatletter
  \newcommand{\xrchap}[1]{%
    \edef\xrc@self{\detokenize{Capitulo#1}}%
    \edef\xrc@job{\jobname}%
    \ifx\xrc@self\xrc@job\else
      \IfFileExists{../#1capitulo/Capitulo#1.aux}%
        {\externaldocument{../#1capitulo/Capitulo#1}}{}%
    \fi
  }
  \makeatother
  \xrchap{01}\xrchap{02}\xrchap{03}\xrchap{04}
  \xrchap{05}\xrchap{06}\xrchap{07}\xrchap{08}
  \xrchap{09}\xrchap{10}\xrchap{11}\xrchap{12}
  \xrchap{13}
\fi
% \PassOptionsToPackage{gray}{xcolor}

% fig_aliasing
% (a) ws > 2 wM (ws=3, wM=1): copias separadas, con separación entre ellas.
% (b) ws = 2 wM (ws=2, wM=1): caso limite, las copias se tocan exactamente en +-wM,
%     donde el borde derecho de la central y el izquierdo de la vecina coinciden.
% (c) ws < 2 wM (ws=1.6, wM=1): en rojo punteado cada copia por separado; en azul
%     el espectro que queda, que es la SUMA de todas. Con wM=1 y ws=1.6 dos copias
%     vecinas se pisan en una franja de ancho 0.4 y ahí la suma vale 0.34 constante
%     (0.85*(1-x) + 0.85*(x-0.6)), de modo que la envolvente alterna picos de 0.85
%     en cada centro con mesetas planas de 0.34. Las copias se repiten hasta los
%     bordes del eje (+-4.8 = 3 ws), donde quedan cortadas a la mitad.

\begin{document}
\begin{tikzpicture}[>={Latex[length=4pt]}, font=\scriptsize, xscale=0.83,
    tri/.style={utnblue, thick, fill=utnblue!12},
    tricopy/.style={utnblue!70, thick, fill=utnblue!6},
    suma/.style={utnblue, very thick},
    copia/.style={red!75!black, semithick, dash pattern=on 2pt off 1.6pt},
    nyq/.style={black!45, thin, dashed},
    ax/.style={->, thin, black},
]
    % ===================== (a) sin solapamiento =====================
    \begin{scope}[shift={(0,0)}]
        \node at (-5.1,0.6) {(a)};
        \node[anchor=south, font=\scriptsize] at (0,1.40) {$\omega_s > 2\omega_M$};
        \draw[ax] (-4.8,0) -- (5.1,0) node[right] {$\omega$};
        \draw[tricopy] (-4,0) -- (-3,0.85) -- (-2,0) -- cycle;
        \draw[tri]     (-1,0) -- (0,0.85)  -- (1,0)  -- cycle;
        \draw[tricopy] ( 2,0) -- (3,0.85)  -- (4,0)  -- cycle;
        % separación entre la copia central y su vecina
        \draw[<->, black!55] (1,0.30) -- (2,0.30);
        \node[black!60, anchor=south, font=\tiny] at (1.5,0.32) {separación};
        % eje de ordenadas y origen
        \draw[ax] (0,-0.05) -- (0,1.05);
        \node[below, font=\tiny] at (0,-0.05) {$0$};
        % marcas del eje
        \draw[thin] (1,0.05) -- (1,-0.05) node[below, font=\tiny] {$\omega_M$};
        \draw[thin] (2,0.05) -- (2,-0.05);
        \draw[thin] (3,0.05) -- (3,-0.05) node[below, font=\tiny] {$\omega_s$};
        \draw[thin] (-3,0.05) -- (-3,-0.05) node[below, font=\tiny] {$-\omega_s$};
        \node[anchor=north, font=\tiny] at (2,-0.42) {$\omega_s\!-\!\omega_M$};
        \draw[thin, -{Latex[length=3pt]}] (2,-0.42) -- (2,-0.13);
    \end{scope}

    % ===================== (b) caso limite =====================
    \begin{scope}[shift={(0,-2.9)}]
        \node at (-5.1,0.6) {(b)};
        \node[anchor=south, font=\scriptsize] at (0,1.40) {$\omega_s = 2\omega_M$};
        \begin{scope}
            \clip (-4.85,-0.05) rectangle (4.85,1.05);
            \foreach \c in {-4,-2,2,4} {
                \draw[tricopy] ({\c-1},0) -- (\c,0.85) -- ({\c+1},0) -- cycle;
            }
            \draw[tri] (-1,0) -- (0,0.85) -- (1,0) -- cycle;
        \end{scope}
        \draw[ax] (-4.8,0) -- (5.1,0) node[right] {$\omega$};
        % las copias se tocan justo en wM
        \node[black!60, anchor=south, font=\tiny] at (1,0.58) {contacto};
        \draw[black!55, thin, -{Latex[length=3pt]}] (1,0.56) -- (1,0.10);
        % eje de ordenadas y origen
        \draw[ax] (0,-0.05) -- (0,1.05);
        \node[below, font=\tiny] at (0,-0.05) {$0$};
        % marcas del eje: los dos bordes que se comparan caen en el mismo punto
        \draw[thin] (1,0.05) -- (1,-0.05);
        \draw[thin] (2,0.05) -- (2,-0.05) node[below, font=\tiny] {$\omega_s$};
        \draw[thin] (-2,0.05) -- (-2,-0.05) node[below, font=\tiny] {$-\omega_s$};
        \node[anchor=north, font=\tiny] at (1,-0.42) {$\omega_M\!=\!\omega_s\!-\!\omega_M$};
        \draw[thin, -{Latex[length=3pt]}] (1,-0.42) -- (1,-0.13);
    \end{scope}

    % ===================== (c) con solapamiento =====================
    \begin{scope}[shift={(0,-5.8)}]
        \node at (-5.1,0.6) {(c)};
        \node[anchor=south, font=\scriptsize] at (0,1.40) {$\omega_s < 2\omega_M$};
        \begin{scope}
            \clip (-4.8,-0.05) rectangle (4.8,1.05);
            % relleno bajo la suma
            \fill[utnblue!12]
                (-4.8,0.85) -- (-4.2,0.34) -- (-3.8,0.34) -- (-3.2,0.85)
                -- (-2.6,0.34) -- (-2.2,0.34) -- (-1.6,0.85)
                -- (-1.0,0.34) -- (-0.6,0.34) -- ( 0.0,0.85)
                -- ( 0.6,0.34) -- ( 1.0,0.34) -- ( 1.6,0.85)
                -- ( 2.2,0.34) -- ( 2.6,0.34) -- ( 3.2,0.85)
                -- ( 3.8,0.34) -- ( 4.2,0.34) -- ( 4.8,0.85)
                -- ( 4.8,0) -- (-4.8,0) -- cycle;
            % contorno de la suma
            \draw[suma]
                (-4.8,0.85) -- (-4.2,0.34) -- (-3.8,0.34) -- (-3.2,0.85)
                -- (-2.6,0.34) -- (-2.2,0.34) -- (-1.6,0.85)
                -- (-1.0,0.34) -- (-0.6,0.34) -- ( 0.0,0.85)
                -- ( 0.6,0.34) -- ( 1.0,0.34) -- ( 1.6,0.85)
                -- ( 2.2,0.34) -- ( 2.6,0.34) -- ( 3.2,0.85)
                -- ( 3.8,0.34) -- ( 4.2,0.34) -- ( 4.8,0.85);
            % cada copia por separado (lo que habría si no se sumaran)
            \foreach \c in {-4.8,-3.2,-1.6,0,1.6,3.2,4.8} {
                \draw[copia] ({\c-1},0) -- (\c,0.85) -- ({\c+1},0);
            }
        \end{scope}
        \draw[ax] (-4.8,0) -- (5.1,0) node[right] {$\omega$};
        % eje de ordenadas y origen
        \draw[ax] (0,-0.05) -- (0,1.05);
        \node[below, font=\tiny] at (0,-0.05) {$0$};
        % linea de Nyquist en +- ws/2 = 0.8
        \draw[nyq] (0.8,-0.05) -- (0.8,0.95);
        \draw[nyq] (-0.8,-0.05) -- (-0.8,0.95);
        \node[black!60, anchor=south, font=\tiny] at (0.8,0.95) {$\tfrac{\omega_s}{2}$};
        % multiplos de ws: el centro de cada copia
        \foreach \x/\lab in {1.6/{\omega_s}, 3.2/{2\omega_s}, 4.8/{3\omega_s},
                             -1.6/{-\omega_s}, -3.2/{-2\omega_s}, -4.8/{-3\omega_s}} {
            \draw[thin] (\x,0.05) -- (\x,-0.05) node[below, font=\tiny] {$\lab$};
        }
        % bordes de la zona de solapamiento: empieza en ws-wM y termina en wM
        \draw[thin] (1,0.05) -- (1,-0.05) node[below, font=\tiny] {$\omega_M$};
        \draw[thin] (0.6,0.05) -- (0.6,-0.05);
        \node[anchor=north, font=\tiny] at (0.6,-0.42) {$\omega_s\!-\!\omega_M$};
        \draw[thin, -{Latex[length=3pt]}] (0.6,-0.42) -- (0.6,-0.13);
        % etiquetas de los dos trazos
        \node[utnblue, anchor=south, font=\tiny] at (4.0,0.92) {suma};
        \draw[utnblue, thin, -{Latex[length=3pt]}] (4.0,0.90) -- (4.0,0.40);
        \node[red!75!black, anchor=south, font=\tiny] at (-4.0,0.92) {copias};
        \draw[red!75!black, thin, -{Latex[length=3pt]}] (-4.0,0.90) -- (-4.0,0.24);
    \end{scope}
\end{tikzpicture}
\end{document}
